\section{The canonical Schur resolution}
\label{app:schur-resolution}

We prove Theorem~\ref{thm:full-schur-resolution}. Retain
$B=K\oplus\mathfrak m$, $S=S_B$, and the shapes $\rho_i$, degrees $d_i$
and differences $e_i=\mu_i-\mu_{i+1}+1$ from
Section~\ref{sec:schur-resolution-statements}. The essential point is a
natural contraction for square-zero multiplication. Its count filtration
makes the perturbation for an arbitrary product finite.

\subsection{A dimension-independent contraction}

Put $E=K1\oplus\mathfrak m$ and give
$B_0=K\oplus\mathfrak m$ the square-zero product. On
\[
                        C_i^0=\Lambda^i\mathfrak m\otimes E^{\otimes n}
\]
let $d_0$ be its moment differential: remove a wedge input and place
it in each empty coefficient position. Define $c:C_i^0\to C_{i+1}^0$
by removing each occupied coefficient $x$ in turn, replacing it by
$1$, and sending the exterior factor $\omega$ to $x\wedge\omega$.
These formulas are natural for every linear map of $\mathfrak m$;
they require no bilinear form. Both operators preserve the count
\[
           k=\text{number of occupied positions}+\text{exterior degree}.
\]
They satisfy $d_0^2=c^2=0$. With any temporary real orthonormal basis,
$c$ is the adjoint of $d_0$ for the standard tensor and exterior
inner products.

\begin{lemma}[The content Laplacian]
\label{lem:schur-content-laplacian}
Let $\cont(\alpha)=\sum_{(r,s)\in\alpha}(s-r)$. On label
$\mu\vdash n$ and polynomial Schur type $\lambda\vdash k$ in $C_i^0$,
the operator $L=d_0c+cd_0$ is the scalar
\begin{equation}\label{eq:schur-content-laplacian}
           \cont(\mu)-\cont(\lambda)+n-\binom{n-k+1}{2}.
\end{equation}
Here the binomial denotes its polynomial value, also for negative
arguments. There is a natural contraction onto $\mathcal H=\ker L$,
with inclusion $j$, projection $p$, and homotopy $h$ satisfying
\[
 d_0h+hd_0=\id-jp,\qquad h^2=0,\qquad hj=ph=0.
\]
On label $\mu$, the harmonic space in homological degree $i$ is
$\mathcal H_i\cong\Schur_{\rho_i}\mathfrak m$, in count $d_i$.
\end{lemma}

\begin{proof}
First let $q=\dim\mathfrak m<\infty$. On the coefficient module
$\Schur_\mu(E)$, write $E_{ab}$ for the matrix-unit action, with
index $0$ denoting the unit. Exterior creation and contraction give
\[
 d_0=\sum_{a=1}^q E_{a0}\iota_a,
 \qquad c=\sum_{a=1}^qE_{0a}\varepsilon_a,
 \qquad
 L=\sum_aE_{a0}E_{0a}+iE_{00}-\sum_{a,b}E_{ab}\varepsilon_b\iota_a.
\]
Let $\Omega_q=\sum_{a,b}E_{ab}E_{ba}$ be the quadratic Casimir.
Its eigenvalue on $\Schur_\lambda K^q$ is
$2\cont(\lambda)+q|\lambda|$. The cross-Casimir identity rewrites
the preceding formula as
\[
 2L=\Omega_{q+1}^{\mu}-\Omega_q^{\mathrm{total}}
       +\Omega_q^{\Lambda^i}
       -E_{00}^2+(2i-q-1)E_{00}+n.
\]
Now $E_{00}=n-k+i$ and
$\Omega_q^{\Lambda^i}=i(q+1-i)$. Substitution cancels $q$ and $i$
and gives \eqref{eq:schur-content-laplacian}.

Over the real numbers the adjoint identity implies
$\ker L=\ker d_0\cap\ker c$. Let $L^+$ invert every nonzero
eigenvalue and be zero on the kernel. Then $h=cL^+$ and the spectral
projection $jp$ give the asserted contraction. Its formulas are
rational: the eigenvalues in \eqref{eq:schur-content-laplacian}
are integers and the Schur projections are natural. They therefore
hold over every characteristic-zero field and commute with arbitrary
linear maps, including maps that lower dimension.

To identify the harmonic spaces, apply the pure-resolution theorem
\cite[Theorem~2.2]{SamWeyman} to the square-zero presentation, with
$\alpha=\tau$ and $\beta=\tau^+$. Its successive generator shapes
are precisely $\rho_i$, and its degree shifts, including the initial
shift $|\tau|$, are $d_i$. Hence the square-zero Koszul homology is
$\Schur_{\rho_i}\mathfrak m$ in count $d_i$. This harmonic copy
occurs once: its coefficient branching shape deletes row $i+1$
of $\mu$, and its exterior Pieri strip adds one box to each of the
first $i$ rows. The case $i=0$ is minimum occupancy.

For infinite-dimensional $\mathfrak m$, each fixed-count construction
passes through finite-dimensional vector subspaces. The scalar formula
is independent of their dimension, so the contractions are compatible
under inclusion and define the same natural operators on their union.
\end{proof}

The diagram convention in \cite{SamWeyman} is transposed, while its
highest-weight tuples agree with those used here. This argument uses
its full pure-resolution theorem; the earlier presentation used only
the Pieri composition lemma.

\subsection{Multiplication as a finite perturbation}

\begin{proof}[Proof of Theorem~\ref{thm:full-schur-resolution}]
We first transfer the collision differential to compute moment homology.
We then adjoin formal trace variables and transfer once more to obtain
the free resolution; setting those variables to zero recovers the first
complex. Both transfers terminate because they lower coefficient count.

For the original algebra, write the moment differential as
$d_B=d_0+\delta$. The collision operator $\delta$ removes a wedge
input and multiplies it into each occupied position. It lowers count
by one, while $h$ preserves count. The perturbation lemma, with our
convention $d_0h+hd_0=\id-jp$, gives the transferred differential
\[
 p\delta(1+h\delta)^{-1}j
       =\sum_{r\ge0}(-1)^r p\delta(h\delta)^rj.
\]
Every series is finite on each vector. This is the standard formula
of \cite[\S2, especially~2.4]{Crainic}, with its homotopy replaced by
$-h$. A term from $\mathcal H_i$ to $\mathcal H_{i-1}$ must lower
count from $d_i$ to $d_{i-1}$, so only one term survives:
\begin{equation}\label{eq:schur-small-differential}
 \partial_i=(-1)^{e_i-1}
          p_{i-1}\delta(h\delta)^{e_i-1}j_i.
\end{equation}
Thus the small complex computes all moment homology, with no choices
of higher splittings.

To obtain an actual free resolution, use
$C_{\mathrm{free}}=S\otimes\Lambda^\bullet\mathfrak m
                    \otimes\Schur_\mu(E)$.
Let $d_T$ remove a wedge input $b$ and multiply the $S$ factor by
$T_b$. Subtracting $d_T$ enforces the action $T_b=p_n(b)$.
The differential $d_B-d_T$ resolves $\Schur_\mu(B)$ under
$f(T)\otimes v\mapsto f(p_n)v$. Indeed, these are polynomial-variable
elimination operators $p_n(b)-T_b$; eliminating any finite set is
exact by successive substitution, and their directed union handles
infinitely many variables.

Extend the preceding contraction $S$-linearly, starting with $d_0$,
and perturb it by $\mathcal A=\delta-d_T$. Both terms lower
\emph{coefficient} count by one; the increase of formal polynomial
degree under $d_T$ is irrelevant to this count. Therefore the
transferred free differential is the finite formula
\begin{equation}\label{eq:schur-free-differential}
 D_i=(-1)^{e_i-1}
       p_{i-1}\mathcal A(h\mathcal A)^{e_i-1}j_i.
\end{equation}
It gives a chain homotopy equivalent free complex with the stated
terms, and hence a resolution. Setting $T=0$ gives
\eqref{eq:schur-small-differential}.

For the resolution filtration, give a formal trace count one and
the generator of $F_i$ count $d_i$. A term with no collisions has
formal degree $e_i$ and preserves this total count; every collision
lowers it. At $\delta=0$ the complex is a minimal equivariant pure
resolution, so its leading maps are the classical nonzero Pieri maps.

The corrected harmonic inclusion is $(1+h\mathcal A)^{-1}j$.
On degree zero, $\mathcal A j_0=0$, so its augmentation is exactly
the minimum-occupancy map. Choose the Schur identification line
for $\mathcal H_1$ so that the leading part of $D_1$ equals the
previously fixed Pieri inclusion. Every lower-count term lies below
$d_1$, where the omitted Schur summands are absent. The unique lift
of that leading relation is therefore the graph lift~\eqref{eq:schur-relation-graph}, so
$D_1=R_{\mu,B}$ literally with this normalization. The earlier
quotient map used $\beta_{\mu,B}=-\operatorname{ev}_{T=0}R_{\mu,B}$;
accordingly $\partial_1=-\beta_{\mu,B}$. The kernels and cokernels
are unchanged. For example, before this line normalization the raw
adjoint convention at $n=1$, $\mathfrak m=K a$, $a^2=ca$ gives
$D_1=cT_a-T_a^2$.

All operators are built from natural Schur projections, unit and
augmentation, and multiplication. This proves the claimed functoriality
and compatibility with gradings. It gives algebra-automorphism
equivariance for general $B$; full $\operatorname{GL}(\mathfrak m)$
equivariance belongs to the square-zero specialization.
\end{proof}

